% Render: python3 /Users/junix/projects/skills/create-tikz-infographic/scripts/render_tikz.py docs/architecture-infographic.tex --workdir /Users/junix/projects/xcrawl --out-dir /tmp/tikz-xcrawl --engine xelatex --dpi 300 --svg --strict-warnings --strict-geometry --expect-text "xcrawl"
% Claim: 调用者经入口进入核心与支撑模块，最终得到可观察结果。
% Reader question: xcrawl 的系统边界、主路径与验证边界是什么？
\documentclass[border=12pt,varwidth=17cm]{standalone}
\usepackage[UTF8,fontset=none]{ctex}
\IfFontExistsTF{Source Han Serif SC}{
  \setmainfont[BoldFont={Source Han Serif SC Bold}]{Source Han Serif SC}
  \setCJKmainfont[BoldFont={Source Han Serif SC Bold}]{Source Han Serif SC}
}{\setCJKmainfont[BoldFont={FandolSong-Bold}]{FandolSong-Regular}}
\IfFontExistsTF{Sarasa Mono SC}{\setmonofont{Sarasa Mono SC}}{\setmonofont{Latin Modern Mono}}
\usepackage{xcolor}
\usepackage{tikz}
\usetikzlibrary{arrows.meta,backgrounds,calc,fit,positioning}
\pgfdeclarelayer{edge layer}
\pgfsetlayers{background,edge layer,main}
\definecolor{Paper}{HTML}{F8FAFC}
\definecolor{Ink}{HTML}{17212B}
\definecolor{Slate}{HTML}{5E6B7A}
\definecolor{Fog}{HTML}{D7DCE3}
\definecolor{Ocean}{HTML}{2563EB}
\definecolor{OceanPale}{HTML}{E8F0FE}
\definecolor{External}{HTML}{EEF1F5}
\tikzset{
  every node/.style={font=\rmfamily},
  component/.style={draw=Ocean,fill=OceanPale,rounded corners=2mm,line width=.8pt,minimum width=30mm,minimum height=12mm,align=center,font=\footnotesize\bfseries,text=Ink},
  external/.style={component,draw=Slate!70,fill=External},
  note/.style={font=\footnotesize,text=Slate,align=center},
  group/.style={draw=Fog,fill=white,rounded corners=3mm,inner sep=5mm},
  group title/.style={font=\footnotesize\bfseries,text=Slate},
  flow/.style={-{Latex[length=2.4mm,width=1.7mm]},draw=Ocean,line width=1.05pt,rounded corners=2mm},
  inferred/.style={-{Latex[length=2.1mm,width=1.5mm]},draw=Slate!75,densely dotted,line width=.8pt,rounded corners=2mm},
  edge label/.style={fill=Paper,inner sep=1.5pt,font=\footnotesize,text=Slate}
}
\begin{document}
\begin{tikzpicture}[x=1cm,y=1cm]
  \begin{pgfonlayer}{background}\fill[Paper,rounded corners=2mm] (-.5,-.55) rectangle (16.5,9.35);\end{pgfonlayer}
  \node[anchor=west,font=\footnotesize\bfseries,text=Ocean] at (0,8.88) {项目架构信息图 · rust};
  \node[anchor=west,text width=15.8cm,align=left,font=\Large\bfseries,text=Ink] at (0,8.25) {xcrawl};
  \node[anchor=west,font=\small,text=Slate] at (0,7.42) {rust 组件：入口、模块、依赖与验证边界};
  \node[external] (caller) at (1.55,5.15) {调用者\\{\footnotesize 人或上游系统}};
  \node[component] (interface) at (5.05,5.15) {CLI\\{\footnotesize 主程序入口}};
  \node[component,line width=1.2pt] (core) at (8.55,5.15) {crawler\\{\footnotesize 核心职责}};
  \node[component] (output) at (14.45,5.15) {命令结果};
  \node[component] (support) at (6.75,2.75) {fetch\\{\footnotesize 支撑职责}};
  \node[external] (dependency) at (10.35,2.75) {async trait\\{\footnotesize 包清单边界}};
  \node[external] (verify) at (13.95,2.75) {测试 / 示例\\{\footnotesize 行为复核}};
  \begin{scope}[on background layer]
    \node[group,fit=(interface)(core)(support),inner xsep=5mm] (system) {};
  \end{scope}
  \node[group title,anchor=south west] at ($(system.north west)+(1mm,2mm)$) {仓库内系统};
  \begin{pgfonlayer}{edge layer}
    \draw[flow] (caller.east) -- node[edge label,above] {调用} (interface.west);
    \draw[flow] (interface.east) -- node[edge label,above] {分发} (core.west);
    \draw[flow] (core.east) -- node[edge label,above] {返回} (output.west);
    \draw[inferred] (interface.south) -- ++(0,-5mm) -| node[edge label,near end,below] {协作?} (support.west);
    \draw[inferred] (core.south) -- ++(0,-5mm) -| node[edge label,near end,below] {依赖?} (dependency.west);
    \draw[inferred] (output.south) -- node[edge label,right] {校验?} (verify.north);
  \end{pgfonlayer}
  \node[note,anchor=west,text width=15.8cm,align=left] at (0,.72) {主路径：调用者 $\rightarrow$ CLI $\rightarrow$ crawler $\rightarrow$ 命令结果。点线表示由入口与源码布局推得的关系，不表达性能或调用次数。};
  \draw[flow] (0,.22) -- (.9,.22);\node[note,anchor=west] at (1.05,.22) {确定边界};
  \draw[inferred] (3.25,.22) -- (4.15,.22);\node[note,anchor=west] at (4.3,.22) {推断关系};
  \node[note,anchor=east] at (16,.22) {代码快照 135d27e};
\end{tikzpicture}
\end{document}
